\documentclass[11pt]{article}
\usepackage[T1]{fontenc}
\usepackage{lmodern,amsmath,amssymb,amsthm}
\usepackage[margin=1in]{geometry}
\usepackage{microtype}
\usepackage[hidelinks]{hyperref}
\usepackage{enumitem}
\setlist{nosep}
\setlength{\emergencystretch}{2em}
\newtheorem{theorem}{Theorem}[section]
\newtheorem{lemma}[theorem]{Lemma}
\newtheorem{corollary}[theorem]{Corollary}
\newcommand{\D}{\mathcal D}
\newcommand{\oct}{\operatorname{oct}}
\title{A Triangle-Deletion Criterion for Tuza's Conjecture\\via Equality in Ryser's Bound}
\author{AI Math Experiments: TUZA-01}
\date{28 October 2026}
\begin{document}
\maketitle
\begin{abstract}
Let $G$ be a finite simple graph and $T$ a triangle. We show that
$\tau(G)\le2\nu(G)$ whenever the triangle hypergraph of $G-E(T)$ is
3-partite. The proof combines Aharoni's tripartite bound with the
Haxell--Narins--Szab\'o characterization of equality. Its graphical
consequence supplies a maximum packing whose triangles avoid a prescribed
independent vertex set. In particular, Tuza's inequality holds when deleting
at most four vertices makes $G$ bipartite. Adapting the proof of Liu--Zeng
then gives, for every fixed real $c<5/2$, the inequality for all sufficiently
large $n$-vertex graphs with $e(G)\ge n^2/4-cn$.
\end{abstract}
\noindent\textbf{Status.} This is an AI-assisted, unrefereed proof manuscript.
No external specialist review has occurred. Originality is unconfirmed:
the searches recorded with this archive found no exact prior equivalent.
The ordinary mathematical proof uses published deep theorems. The accompanying
Lean development checks the new deduction conditionally on named
\texttt{AharoniInput} and \texttt{HNSInput} hypotheses; it does not formalize
those theorems or the density corollary. The general Tuza conjecture remains open.

\section{Definitions, statement, and context}
All graphs are finite, simple, and undirected. A triangle is a three-vertex
clique, identified with its three edges when discussing packing and covering.
For a triangle family $\mathcal F$, let $\nu(\mathcal F)$ be the maximum
number of pairwise edge-disjoint members, and let $\tau(\mathcal F)$ be the
minimum number of graph edges meeting every member. Both parameters are zero
for the empty family. Write $\nu(G),\tau(G)$ for the full triangle family.
The odd-cycle transversal number $\oct(G)$ is the minimum number of
\emph{vertices} whose deletion leaves a bipartite graph.

The triangle hypergraph $\D(G)$ has vertex set $E(G)$ and a hyperedge for
the edge set of every triangle. Its matching and vertex-cover numbers are
$\nu(G)$ and $\tau(G)$. It is linear: distinct triangles share at most one
edge. It is 3-partite precisely when the graph edges can be labeled with
three labels so that every triangle receives all three. Labels on other
incident edges need not differ. This condition is often called
triangle-3-colorability and is broader than ordinary four-colorability.

\begin{theorem}[Triangle-deletion criterion]\label{thm:main}
Let $T$ be a triangle of $G$ and put $H=G-E(T)$, retaining all vertices.
If $\D(H)$ is 3-partite, then $\tau(G)\le2\nu(G)$.
\end{theorem}

Tuza's inequality for triangle-3-colorable graphs, including four-colorable
graphs, is known from Lakshmanan--Bujt\'as--Tuza \cite{LBT} and follows from
Aharoni's theorem \cite{Aharoni}. The new deduction considered here is the
addition of the three edges of a triangle to such a remainder. Merely
adding these three edges to a cover loses one unit in the critical case.
Equality structure repairs that loss. We use only the specific consequence
of Haxell--Narins--Szab\'o (HNS) \cite{HNS} proved below, rather than a
classification of all graphical equality cases.

There are antecedents for $K_4$ equality arguments in planar graphs
\cite{CHM} and for fractional triangle covers \cite{HKT}. The latter concerns
$\tau^*$, and integer equality $\tau=2\nu$ cannot simply be substituted for
fractional equality. The elementary intersecting-family lemma below is also
known \cite{AZ}. No novelty is claimed for these inputs. Liu--Zeng's
preprint \cite{LZ} proves the odd-cycle-transversal-three case and the
fixed-deficit density range $c<2$; our corollaries concern four vertices and
$c<5/2$, respectively.

\section{The equality consequence of the published inputs}
\paragraph{Aharoni's bound.}
Every finite 3-partite 3-uniform hypergraph $\mathcal H$ satisfies
$\tau(\mathcal H)\le2\nu(\mathcal H)$ \cite{Aharoni}.
An accessible formulation and derivation also appear in HNS, Section 3,
Proposition 3.2(iii) and the following paragraph (author-manuscript p.~9).
Thus $\D(J)$ 3-partite implies $\tau(J)\le2\nu(J)$.

\paragraph{HNS's equality characterization.}
HNS Theorem 1.1 characterizes 3-partite 3-uniform hypergraphs with
$\tau=2\nu$ as home-base hypergraphs \cite{HNS}. We use the 18 March 2018
author manuscript: Definitions 7.1--7.4, pp.~19--20, and Proposition 7.5,
p.~21. The earlier arXiv version has different numbering. HNS permits
parallel hyperedges; simple hypergraphs are included.

Here are exactly the features of a home-base partition needed for extraction.
The partition is denoted $(\mathcal A,\mathcal R,W)$: the disjoint pieces
in $\mathcal A$, the disjoint triples in $\mathcal R$, and $W$ partition
the hypervertices, and $|\mathcal A|+|\mathcal R|=\nu(\mathcal H)$.
Each $\mathcal A$-piece is a truncated multi-Fano configuration on six
vertices, containing four distinct underlying triples, any three of which
form a loose 3-cycle. For a triple $R=\{r_1,r_2,r_3\}\in\mathcal R$,
with $r_i$ in part $V_i$, each auxiliary bipartite graph $B_i$ has a
matching saturating $\mathcal R$. A neighbor $w\in W\cap V_i$ of $R$
means that $(R\setminus\{r_i\})\cup\{w\}$ is a hyperedge.
These are consequences of the cited definitions, not additional hypotheses.

\begin{lemma}[Disjoint loose cycles]\label{lem:cycles}
If $\mathcal H$ is 3-partite and 3-uniform with
$\tau(\mathcal H)=2\nu(\mathcal H)=2h$, then it contains $h$ loose
3-cycles with pairwise disjoint vertex supports. A loose 3-cycle here consists
of three hyperedges on six vertices whose pairwise intersections are three
distinct singletons.
\end{lemma}
\begin{proof}
For $h=0$, take the empty collection. Otherwise use the HNS partition.
Choose a saturating matching in each $B_i$, and write $w_i(R)$ for the
vertex matched to $R$. For each $R$ the hyperedges
\begin{equation}\label{eq:cycles}
\{r_2,r_3,w_1(R)\},\quad
\{r_1,r_3,w_2(R)\},\quad
\{r_1,r_2,w_3(R)\}
\end{equation}
form the required loose cycle. Its six vertices are distinct since
$W$ is disjoint from the central triples and the three $w_i(R)$ lie in
different parts. Across different $R$'s, the same-part chosen vertices
differ by the matching property; different-part vertices differ by the
tripartition. None belongs to any central triple or $\mathcal A$-piece.
Thus these supports are disjoint simultaneously; no compatibility between
the three matchings is missing. Choose also three distinct underlying
hyperedges in each $\mathcal A$-piece. These give disjoint cycles, for a
total of $|\mathcal A|+|\mathcal R|=h$. This is the extraction also used
in the proof of HNS Proposition 7.5. Isolated vertices may remain in $W$.
\end{proof}

\begin{lemma}[Graphical equality consequence]\label{lem:k4}
If $\D(J)$ is 3-partite and $\tau(J)=2\nu(J)=2h$, then $J$ contains
$h$ pairwise edge-disjoint copies of $K_4$.
\end{lemma}
\begin{proof}
Apply Lemma~\ref{lem:cycles} to $\D(J)$. Consider one loose cycle.
Name two graph triangles $uvx$ and $uvy$, where the four vertices are
distinct. The third shares with each an edge other than $uv$.
Its two shared edges lie in $\{ux,vx\}$ and $\{uy,vy\}$, respectively.
They must meet, since they belong to one triangle. The cross choices
$ux,vy$ and $vx,uy$ are disjoint, so the third triangle is $uxy$ or $vxy$.
The union therefore contains exactly the six edges
$uv,ux,vx,uy,vy,xy$ of a $K_4$. Disjoint hypervertex supports mean disjoint
graph-edge sets. The resulting $K_4$ copies may share graph vertices.
\end{proof}

\begin{lemma}[Faces avoiding an independent set]\label{lem:faces}
Under the hypotheses of Lemma~\ref{lem:k4}, for every independent vertex
set $U$ in $J$ there is a maximum triangle packing all of whose vertices
lie outside $U$.
\end{lemma}
\begin{proof}
Each $K_4$ meets $U$ in at most one vertex. Choose its opposite face if
it meets $U$, and any face otherwise. The chosen faces are edge-disjoint,
avoid $U$, and number $h=\nu(J)$. For $h=0$ use the empty packing.
\end{proof}

\section{The triangle-deletion proof and the OCT corollary}
\begin{lemma}[Two-edge cover]\label{lem:intersecting}
Every nonempty pairwise edge-intersecting family of graph triangles has
a cover consisting of at most two graph edges. The family need not be
the full triangle family of its ambient graph.
\end{lemma}
\begin{proof}
A single triangle needs one edge. Otherwise choose distinct members
$abc,abd$ sharing $ab$. Any member avoiding $ab$ must share an edge from
$\{ac,bc\}$ and one from $\{ad,bd\}$. The disjoint cross pairs cannot
be edges of one triangle; hence it must be $acd$ or $bcd$.
If there is no such member, $ab$ covers the family. Otherwise $cd$ is
an actual graph edge and $\{ab,cd\}$ covers the family.
This also gives the $k=3$ case of the known intersecting-family result
in Aharoni--Zerbib, Propositions 3.2 and 3.7 \cite{AZ}.
\end{proof}

\begin{proof}[Proof of Theorem~\ref{thm:main}]
Put $h=\nu(H)$ and $U=V(T)$. A maximum packing in $H$ together with
$T$ gives $\nu(G)\ge h+1$. Aharoni gives $\tau(H)\le2h$.
If the inequality is strict, integrality and a cover augmented by $E(T)$ give
\begin{equation}\label{eq:strict}
\tau(G)\le\tau(H)+3\le2h+2\le2\nu(G).
\end{equation}
Suppose instead $\tau(H)=2h$. The set $U$ is independent in $H$, since
all three possible edges on it have been removed. Lemma~\ref{lem:faces}
gives a packing $\mathcal P$ of size $h$ avoiding $U$.
Let $\mathcal F$ be the triangles of $G$ absent from $H$, and put
$r=\nu(\mathcal F)$. Every member of $\mathcal F$ contains an edge of
$T$, hence at least two vertices of $U$. Consequently every edge of
every member touches $U$. Every edge of $\mathcal P$ has both endpoints
outside $U$. The single fixed packing $\mathcal P$ is therefore
edge-disjoint from \emph{all} members of $\mathcal F$ simultaneously, and
\begin{equation}\label{eq:simultaneous}
\nu(G)\ge h+r.
\end{equation}
Since $T\in\mathcal F$, we have $r\ge1$. If $r=1$, any two members
intersect in an edge, so Lemma~\ref{lem:intersecting} gives
$\tau(\mathcal F)\le2=2r$. If $r\ge2$, the three edges of $T$ give
$\tau(\mathcal F)\le3\le2r$. Combining covers of $H$ and $\mathcal F$ yields
\begin{equation}\label{eq:equality}
\tau(G)\le\tau(H)+\tau(\mathcal F)\le2h+2r\le2\nu(G).
\end{equation}
This includes $h=0$, using an empty packing and no HNS invocation. A graph
with no triangle cannot satisfy the stated premise on $T$.
\end{proof}

For clarity, the choice of packing matters: an arbitrary maximum packing
of $H$ need not avoid the new triangles. For example, take a $K_4$ on
$a,d,e,f$, add the edge $bd$, and then add $T=abc$. The old maximum
packing $\{ade\}$ meets the new triangle $abd$, whereas the face $def$
avoids every new triangle. Lemma~\ref{lem:faces} supplies the needed choice.

\begin{lemma}[Four colors give three edge labels]\label{lem:colors}
If $J$ is four-colorable, then $\D(J)$ is 3-partite.
\end{lemma}
\begin{proof}
Identify vertex colors with $\mathbb F_2^2$ and label an edge by the
sum of its endpoint colors. Properness makes this one of the three
nonzero vectors. A triangle has three distinct vertex colors, so its
three edge labels are distinct. Empty edge classes cause no difficulty.
\end{proof}

\begin{corollary}\label{cor:oct}
If $\oct(G)\le4$, then $\tau(G)\le2\nu(G)$.
\end{corollary}
\begin{proof}
Choose $S\subseteq V(G)$ with $|S|\le4$ and $G-S$ bipartite. If
$G[S]$ is bipartite, use two colors on $S$ and two different colors on
its complement. Apply Lemma~\ref{lem:colors} and Aharoni directly.
Otherwise $G[S]$ contains an odd cycle, which on at most four vertices
must be a triangle $T$. After deleting $E(T)$ the only remaining edges
inside $S$ join the possible fourth vertex to $V(T)$, so they form a
subgraph of a star. Both $H[S]$ and $H-S$ are bipartite. Thus $H$ is
four-colorable, and Lemma~\ref{lem:colors} and Theorem~\ref{thm:main} apply.
The argument includes all smaller transversals and all configurations on $S$.
\end{proof}

\section{A fixed-deficit density corollary}
The following is a derived corollary using the proof of Liu--Zeng
\cite{LZ}, not a theorem quoted from their preprint. We give the modified
estimates and the small-packing argument in detail, identifying the
unchanged inputs. This section has not been formalized in Lean.

\begin{corollary}[Adaptation of Liu--Zeng]\label{cor:density}
For every fixed real $c<5/2$ there exists $N(c)$ such that, for every
integer $n\ge N(c)$ and every $n$-vertex graph $G$,
$e(G)\ge n^2/4-cn$ implies $\tau(G)\le2\nu(G)$.
\end{corollary}

\paragraph{Imported approximation and elementary estimates.}
Let $\nu^*(G)$ be the fractional triangle-packing optimum. We use
$\nu^*(G)-\nu(G)=o(n^2)$, due to Haxell--R\"odl \cite{HR}; alternatively
Yuster's Theorem 1.2 \cite{Yuster}, with graph family $\{K_3\}$, suffices.
Removing isolates leaves both optima unchanged. If the remaining order
is large, its $o(m^2)$ error is bounded by the desired error in $n^2$;
if bounded, its whole fractional optimum is a fixed constant. This
handles Yuster's convention excluding isolated vertices.
Finite linear-programming duality identifies $\nu^*$ with the fractional
edge-cover optimum.

We also use Liu--Zeng Lemma 2.4: if $S\subseteq E(J)$ has at most one
edge per triangle, $C=E(J)\setminus S$, and $|C|\ge\beta n^2$ for
$0<\beta<1/2$, then
\begin{equation}\label{eq:saving}
\tau(J)\le |C|-\beta^3n^2/36.
\end{equation}
This lemma has no density-deficit parameter. Their proof deletes a
suitable cut complement, starting from
$\tau(J)\le|C|-e_C(A,B)+e_S(A)+e_S(B)$ for disjoint $A,B$.
Its independent-neighborhood partition gives savings at least
$\beta^3n^2/6$, $\beta^3n^2/36$, or $\beta n^2/24$ according to the
location of a large class of $C$-edges. Thus the common constant in
\eqref{eq:saving} remains unchanged here.

\begin{proof}[Proof of Corollary~\ref{cor:density}]
It suffices to treat $0\le c<5/2$, since a negative $c$ has a stronger
hypothesis than $c=0$. Put $k=\nu(G)$, $t=\tau(G)$ and take $n$
sufficiently large, as specified below.

\paragraph{Large packing: $k\ge10^{-5}n^2$.}
This is Liu--Zeng's Section 3, Case 1, unchanged: it uses no edge-count
assumption or bound on $c$. Here is its quantitative core. If $t>2k$,
choose the approximation error at most $10^{-18}n^2$ and write
$d=\nu^*(G)-t/2$. The fractional-cover threshold argument in their
equations (2)--(6), applied to weights in decreasing order, gives
$0\le d<10^{-18}n^2$. Partition the edges into $L,M,F$ according as
their weights lie in $[0,1/6)$, $[1/6,2/3)$, $[2/3,1]$. The same
argument gives
\begin{equation}\label{eq:large}
|F|+|M|-t\le12d,\qquad |F|\le6d,\qquad |M|>10^{-5}n^2.
\end{equation}
For precision, the threshold inequalities are
$2f_i+f_{2t+1-i}\ge1$ for $1\le i\le t$; summing bounds the weight
after the first $t$ edges by $2d$, and clipping weights at $1/2$ bounds
their total excess by $d$. These yield \eqref{eq:large} exactly as in
the cited case. A triangle in $G-F$ cannot have two $L$-edges, since
their weights together with the third weight would sum to less than one.
Apply \eqref{eq:saving} with $\beta=10^{-5}$ to obtain
\begin{equation}\label{eq:contradiction}
\begin{aligned}
t&\le |F|+|M|-10^{-15}n^2/36\\
 &<t+(12\cdot10^{-18}-10^{-15}/36)n^2<t,
\end{aligned}
\end{equation}
a contradiction. All inputs to this case are independent of $c$.

\paragraph{A triangle-free decomposition.}
For subsequent cases let $H$ be triangle-free and put
$\delta=n^2/4-e(H)$, where $0\le\delta\le10^{-4}n^2$.
Take a maximum-degree vertex of degree $q$, set $V_1=N_H(v)$ and
$V_2=V(H)\setminus V_1$. Since $V_1$ is independent,
$e(H)+e_H(V_2)\le q(n-q)\le n^2/4$. The number of missing cross
pairs is at most $2\delta$, and each part differs from $n/2$ by at
most $\sqrt{2\delta}$. Remove into $A$ every vertex missing more
than $n/5$ neighbors in the other part; put $B=V_1\setminus A$,
$C=V_2\setminus A$. Counting missing pairs twice gives
\begin{equation}\label{eq:decomposition}
|A|\le20\delta/n\le n/500,\quad |B|,|C|>9n/20.
\end{equation}
Every distinct pair in $B$ has more than $n/20$ common $H$-neighbors
in $C$, and conversely. Hence $B,C$ are independent in $H$ and the
missing $B,C$-pairs number at most $2\delta$. These are Liu--Zeng's
equations (7)--(13), with the deficit retained explicitly.

\paragraph{Intermediate packing: $n^{2/3}<k<10^{-5}n^2$.}
Remove a maximum packing $\mathcal P$ to form $H$. Then
\begin{equation}\label{eq:delta}
0\le\delta\le cn+3k\le3n+3k\le10^{-4}n^2
\end{equation}
eventually. The replacement of $2n$ by $3n$ is the first change to the
source. Use the preceding decomposition. For $x\in A$, augment its
$H$-neighbors in $B$ by the two vertices of each packing triangle
$xbb'$; call this set $L_x$. Define $R_x$ analogously on $C$.
There is no $H$-edge from $L_x$ to $R_x$. To see the exchange when
such an edge $bc$ exists, if both $xb,xc$ are $H$-edges there is a
triangle in $H$. If $xbb'$ is packed and $xc$ is in $H$, replace
$xbb'$ by $xbc$ and $bb'z$, choosing a common $H$-neighbor $z\in C$
of $b,b'$ with $z\ne c$. The symmetric case is identical. If
$xbb'$ and $xcc'$ are both packed, replace them by $xbc,bb'z,cc'w$,
where $z\in C\setminus\{c,c'\}$ and
$w\in B\setminus\{b,b'\}$ are appropriate common $H$-neighbors.
The $n/20$ lower bounds permit these choices. The inserted triangles
are edge-disjoint, and their new edges lie in $H$, so none meets a
retained packing triangle. A packed triangle entirely in $B$ or $C$
similarly has a one-to-two exchange using two distinct common neighbors.

Place each $x\in A$ on the side of the smaller of $L_x,R_x$.
Their product counts missing $B,C$-pairs, so its smaller size is at
most $\sqrt{2\delta}$. A packed triangle contributes one internal
edge, plus two if all its vertices lie on one side. The extra
contributions from triangles meeting $A$ once are already counted
by the augmented sets. Let $Q$ be the wholly internal packed triangles
meeting $A$ at least twice. They can be charged to distinct edges on
$A$ outside $H$. Thus the internal $H$-edges on $A$ plus $|Q|$ number
at most $\binom{|A|}{2}$, and these edges plus $2|Q|$ number at most
$|A|^2$. Deleting all internal edges therefore gives
$t\le k+|A|\sqrt{2\delta}+|A|^2$, hence
\begin{equation}\label{eq:middle}
t\le k+\frac{20\sqrt2}{n}\delta^{3/2}+\frac{400}{n^2}\delta^2.
\end{equation}
Using $(a+b)^{3/2}\le\sqrt2(a^{3/2}+b^{3/2})$,
$(a+b)^2\le2(a^2+b^2)$, and \eqref{eq:delta}, we obtain
\begin{equation}\label{eq:margin}
\begin{aligned}
\frac{t-k}{k}\le{}&120\sqrt3\,n^{-1/6}
 +120\sqrt{3\cdot10^{-5}}\\
 &+7200n^{-2/3}+7200\cdot10^{-5}<1
\end{aligned}
\end{equation}
eventually. The two constant terms sum to approximately $0.729267069$.
For example $n\ge10^{18}$ suffices for this estimate and the elementary
size conditions here. This proves $t\le2k$ in the intermediate case.

\paragraph{Small packing: $k\le n^{2/3}$, including equality.}
Choose a spanning subgraph $G'$ with the fewest edges subject to
$e(G)-e(G')=k-k'$, where $k'=\nu(G')$. Such a graph exists since $G$
qualifies, and
\begin{equation}\label{eq:restore}
\tau(G)\le k-k'+\tau(G').
\end{equation}
Every edge of $G'$ has codegree at most $2k'$. Indeed, if $uv$ had more
common neighbors, any packing in $G'-uv$ rules out at most two such
neighbors per packed triangle. One further page $uvw$ extends it,
forcing $\nu(G'-uv)=k'-1$. This would contradict the edge-minimal
choice of $G'$ (Liu--Zeng Lemma 2.5).
If $k'=0$, \eqref{eq:restore} settles the claim. Otherwise remove a
maximum packing from $G'$ to obtain a triangle-free $H$ with
\begin{equation}\label{eq:small-deficit}
e(H)=e(G)-k-2k'\ge n^2/4-cn-3k\ge n^2/4-5n
\end{equation}
eventually. Thus the source's constant $5n$ is unchanged for $c<5/2$.
The decomposition gives $|A|\le100$. Since $n/20>2k'$ eventually,
its common-neighbor property and the codegree bound make $B,C$
independent even in $G'$.

Place each $x\in A$ on the side with fewer $H$-neighbors. Their product
is at most $10n$, since all pairs between these neighbor sets are missing
cross pairs. The resulting parts $X,Y$ satisfy
\begin{equation}\label{eq:small-partition}
\begin{aligned}
\bigl||X|-n/2\bigr|,\ \bigl||Y|-n/2\bigr|&\le\sqrt{10n}+100,\\
e_{G'}(X)+e_{G'}(Y)&\le4950+100\sqrt{10n}+3k'.
\end{aligned}
\end{equation}
The first term $4950$ bounds edges inside $A$; at most
$100\sqrt{10n}$ internal $H$-edges meet its complement, and all
remaining edges removed in forming $H$ number $3k'$.

Let $S_1,S_2$ be minimum vertex covers of the internal graphs on $X,Y$,
and write $s=|S_1|+|S_2|\le100$; the parts of $A$ themselves cover
these graphs. The following finite-witness fact is Liu--Zeng Lemma 2.6.
A graph with a vertex cover of size $a$ has an induced subgraph on at
most $a^2+2a$ vertices with the same minimum vertex-cover number:
retain the cover and up to $a+1$ outside neighbors of each cover
vertex. A cover of the retained graph of size at most $a$ extends
to the whole graph. Otherwise a missed cover vertex with an omitted
neighbor would have $a+1$ retained neighbors, a contradiction.
The assertion also holds for $a=0$.
Apply this to get witnesses $T_1\subseteq X,T_2\subseteq Y$ with their
respective vertex-cover numbers $|S_1|,|S_2|$ and
$|T_1|+|T_2|\le s^2+2s\le10200$.

For $y\in Y\setminus T_2$, its nonneighbors in $T_1$, together with
one endpoint of each edge among its neighbors in $T_1$, cover $G'[T_1]$.
Summing and using the codegree bound gives at least
\begin{equation}\label{eq:missing}
|S_1|(|Y|-|T_2|)-2k'e_{G'}(T_1)
\end{equation}
missing pairs in $T_1\times(Y\setminus T_2)$. The symmetric count
lies in $(X\setminus T_1)\times T_2$, a disjoint set of pairs.
Use
\begin{equation}\label{eq:witness-numbers}
\begin{aligned}
e_{G'}(T_1)+e_{G'}(T_2)&\le\binom{10200}{2}=52,014,900,\\
|S_1||T_2|+|S_2||T_1|&\le1,020,000.
\end{aligned}
\end{equation}
Subtracting both missing-pair counts from $|X||Y|\le n^2/4$ and
adding the internal edges from \eqref{eq:small-partition} yields
\begin{equation}\label{eq:upper}
e(G')\le\frac{n^2}{4}-\frac{s}{2}n
 +200\sqrt{10n}+1,034,950+104,029,803k'.
\end{equation}
The constant terms are $100s+1,020,000+4950\le1,034,950$;
the coefficient of $k'$ is $2\cdot52,014,900+3$.
Compare with $e(G')\ge n^2/4-cn-k+k'$ and use
$k'\le k\le n^{2/3}$ to obtain
\begin{equation}\label{eq:key}
\left(\frac{s}{2}-c\right)n
\le104,029,803n^{2/3}+200\sqrt{10n}+1,034,950.
\end{equation}
For $n\ge1$ the right side is less than $106,000,000n^{2/3}$.
Put $\varepsilon=5/2-c>0$. If $s\ge5$, \eqref{eq:key} is impossible
for $n>(106,000,000/\varepsilon)^3$. Hence $s\le4$.
Deleting $S_1\cup S_2$ makes $G'$ bipartite. Corollary~\ref{cor:oct}
gives $\tau(G')\le2k'$, and \eqref{eq:restore} gives
$\tau(G)\le k+k'\le2k$.

One sufficient $N(c)$ is the maximum of $10^{18}$, the threshold for
the imported fractional approximation with error $10^{-18}n^2$, and
an integer strictly exceeding $(106,000,000/(5/2-c))^3$.
For $c<0$ use $N(0)$. No numerical threshold for the deep approximation
is asserted. The strict positivity of $5/2-c$ is essential to the
last step: this proof establishes neither the endpoint $c=5/2$ nor a
threshold uniform as $c$ approaches it. This completes the adaptation.
\end{proof}

\section{Formal verification, limitations, and provenance}
The accompanying Lean 4.33.1 project uses Mathlib v4.33.1, with its full
revision and every transitive dependency pinned in the lockfile.
Its declarations \texttt{Tuza.triangle\_deletion} and
\texttt{Tuza.oct\_four} quantify over arbitrary finite simple graphs.
They use actual finite packing maxima and cover minima. A support-trimming
lemma and \texttt{Tuza.graph\_min\_cover} ensure a minimum consists of
actual graph edges despite a permissive intermediate set representation.

Both main declarations explicitly take \texttt{AharoniInput}, the
graphical tripartite inequality, and \texttt{HNSInput}, the graphical
equality consequence in Lemma~\ref{lem:k4}. The home-base extraction and
graphical loose-cycle conversion remain inside this external mathematical
boundary. Face selection, simultaneous compatibility, the arbitrary-family
cover lemma, the strict/equality arithmetic, and the OCT reduction are
proved in Lean. The deep input theorems and the density argument are not.
The recorded transitive axioms are only \texttt{propext},
\texttt{Classical.choice}, and \texttt{Quot.sound}. That report does not
discharge parameterized hypotheses. Small exact tests corroborate finite
definitions and examples; they do not establish an infinite theorem.

The archive records separate Phase 1 research, Validation V1, manuscript
criticism, and publication/reproduction criticism. Those internal reviews
are not external specialist review or peer review. Bounded literature
searches, including alternative formulations and related HNS applications,
found no exact predecessor; inaccessible or unindexed material and
specialist knowledge remain possible sources of prior art. We do not
claim priority or a solution of the general conjecture.

This manuscript was prepared with AI assistance under the neutral project
attribution AI Math Experiments. No human author, affiliation, external
reviewer, or unsupported model configuration is asserted. All reused
mathematics is attributed to the cited sources. The repository retains
the proof sources, a complete Markdown exposition, reproducible checks,
and sanitized research findings; generated PDFs and dependency caches
are excluded from its text-only publication tree.

\bibliographystyle{plain}
\bibliography{references}
\end{document}
